% 18 November 2026. Prepared for Edward Baker with OpenAI Codex LLM assistance.
% Actual model and effort unavailable; helper settings inherited.
% E31 author A. Sources: CM Notes32/34, with E29/E30 scope retained.
\section{Forward finite-block statistics from a counted initial slice}
\label{sec:forward-statistics}

Known information at time zero can be propagated forward before a
backward comparison is attempted. The two operations require different
estimates. Forward root repulsion can improve a lower bound for a block
minimum, whereas backward protection needs a lower interaction or forcing
bound. We first establish the forward statements, retaining the complete
exterior in each one. The results of \cite{BakerCM32,BakerCM34}, assessed
in \cite{BakerCMReview29,BakerCMReview30}, are conditional finite estimates;
they do not supply a new time-wide canonical source certificate.

\subsection{What a complete initial census can imply}

\begin{lemma}[Reality from a complete singleton census]
\label{fwd:singleton}
Let $H(0,\overline z)=\overline{H(0,z)}$, and let $\Omega$ be a bounded
conjugation-invariant region with zero-free boundary. Suppose a complete
census of its zeros consists of finitely many cells whose horizontal
projections are pairwise disjoint. If every occupied cell has count one,
including multiplicity, every zero in $\Omega$ is real and simple.
\end{lemma}
\begin{proof}
If $\rho$ is a zero in one cell, its conjugate is also in $\Omega$.
Completeness places $\overline\rho$ in an occupied cell. Since the real
coordinates agree, disjoint horizontal projections place both occurrences
in the same cell. Its count one forces $\rho=\overline\rho$, and also
forces multiplicity one. No replacement of an imaginary interval by its
midpoint is used. The cells themselves need not be conjugation invariant.
\end{proof}

Thus a complete finite census may contain more useful information than
one scalar subsequently evaluated from it. Ordered disjoint real cells
give interval enclosures for all internal adjacent gaps and, away from the
two end edges, for neighboring triples of gaps. The numerical census of
\cite{BakerOCCensus22} illustrates this deduction in
Appendix~\ref{app:compact-example}, at its inherited winding/count scope.
The deduction supplies neither global reality at time zero nor a
positive-time source or root-velocity bound.

\subsection{A local forward interval without global reality}

Fix a bounded conjugation-invariant domain $\Omega$ with a simple,
piecewise smooth boundary $\Gamma$. Assume the complete heat field is
nonzero on a neighborhood of $\Gamma$ for $0\le t\le s$ and has exactly
$N\ge2$ zeros inside, counting multiplicity. All source regularity and
clearance assumptions concern this whole interval. Initially suppose
these $N$ zeros are real and simple. Exterior zeros may be nonreal.

On a simple-root subinterval write the ordered interior labels as
$x_1<\cdots<x_N$ and define the whole-block deflation
\begin{equation}
 R(t,z)=B(t,z)-\sum_{j=1}^{N}\frac1{z-x_j(t)},\qquad
 U_i=\sum_{j\ne i}\frac1{x_i-x_j}+R(t,x_i),\qquad
 x_i'=2U_i.
 \label{fwd:deflation}
\end{equation}
The definition uses the same complete $B=H_z/H$ as
Section~\ref{sec:lattice}. In particular $R$ retains all exterior poles
and the regularizing terms of the complete field. It is analytic inside
$\Gamma$ after the interior poles have been removed; no independent
infinite first-power sum is introduced.

\begin{proposition}[Forward preservation on a cleared finite domain]
\label{fwd:preservation}
Under the preceding hypotheses the $N$ interior zeros remain real and
simple throughout $0\le t\le s$.
\end{proposition}
\begin{proof}
A simple real zero continues as a real zero by conjugation and local
uniqueness. Suppose there were a first loss of simplicity, and isolate
one of the colliding clusters of $k\ge2$ real labels. Its limiting point
is inside $\Omega$: the time-wide boundary clearance prevents escape or
entry. Choose a small fixed circle separating this cluster from every
other zero near the putative event. Removing all its poles gives an
analytic regular part $R_{\rm cl}$, uniformly bounded near the cluster
up to the event. For example, the monic polynomial of the enclosed roots
has analytic symmetric coefficients obtained from the counted contour;
the remaining factor is analytic and nonzero there.

Set $\bar x=k^{-1}\sum_{i=1}^k x_i$ and
$\mathcal V=\sum_{i=1}^k(x_i-\bar x)^2$. Since the centered coordinates
sum to zero, the node law gives
\begin{align}
 \mathcal V'
 &=4\sum_i(x_i-\bar x)
       \left(\sum_{j\ne i}\frac1{x_i-x_j}+R_{\rm cl}(t,x_i)\right)
       \notag\\
 &=2k(k-1)+4\sum_i(x_i-\bar x)R_{\rm cl}(t,x_i).
 \label{fwd:cluster-variance}
\end{align}
For the second equality, pair the two ordered contributions of each
unordered pair: their quotient is
$[(x_i-\bar x)-(x_j-\bar x)]/(x_i-x_j)=1$.
At the proposed collapse the last term tends to zero, so
$\mathcal V'>0$ sufficiently close to the event. A nonnegative
$\mathcal V$ tending to zero from the left cannot have that property.
This excludes every cluster size, including simultaneous disjoint
collisions. The local continuation and the compact boundary clearance
then cover the entire interval.
\end{proof}

For an isolated pair the same sign appears as $D'=2$ at $D=0$.
Proposition~\ref{fwd:preservation} includes higher clusters and does not
assume reality of the complete divisor. It is nevertheless conditional
on the independent time-wide contour certificate. A time-zero census
does not acquire that certificate, spatial buffers, or any numerical
lower bound for the later gaps.

\subsection{Dissipation of common node-force dispersion}

Let $P=I-\mathbf1\mathbf1^{T}/N$ be the orthogonal projection onto
mean-zero vectors. Define a positive weighted graph Laplacian by
\begin{equation}
 (Lv)_i=\sum_{j\ne i}\frac{v_i-v_j}{(x_i-x_j)^2},\qquad
 h_i=(R_t+2U_iR_z)(t,x_i),\qquad
 a=PU,\quad r_U=\lVert a\rVert_2.
 \label{fwd:laplacian}
\end{equation}
The derivatives of $R$ are taken at a fixed spatial argument before
evaluation at $x_i$. All external and block-boundary effects remain in
$h$.

\begin{proposition}[Force dispersion in the two time directions]
\label{fwd:force-identity}
On a common simple-root interval,
\begin{align}
 U_t&=-2LU+h,\notag\\
 \frac{\dd}{\dd t}\frac{r_U^2}{2}
 &=-2\sum_{i<j}\frac{(U_i-U_j)^2}{(x_i-x_j)^2}
       +(PU)\mathbin{\cdot}(Ph).
 \label{fwd:energy}
\end{align}
In backward time $u=t_*-t$, put
$f_i=(R_u-2U_iR_z)(u,x_i)=-h_i(t_*-u)$. Then
\begin{align}
 U_u&=2LU+f,\notag\\
 \frac{\dd}{\dd u}\frac{r_U^2}{2}
 &=2\sum_{i<j}\frac{(U_i-U_j)^2}{(x_i-x_j)^2}
       +(PU)\mathbin{\cdot}(Pf).
 \label{fwd:backward-energy}
\end{align}
\end{proposition}
\begin{proof}
Differentiate each finite reciprocal in \eqref{fwd:deflation}, using
$x_i'=2U_i$:
\[
 \frac{\dd}{\dd t}\frac1{x_i-x_j}
 =-\frac{2(U_i-U_j)}{(x_i-x_j)^2}.
\]
The derivative of $R(t,x_i(t))$ is $h_i$, proving the vector equation.
Since $L\mathbf1=0$, $a\cdot LU=a\cdot La$; pairing $i,j$ gives
$a\cdot La=\sum_{i<j}(U_i-U_j)^2/(x_i-x_j)^2$.
Taking the scalar product with $a$ proves \eqref{fwd:energy}.
Reversing physical time proves \eqref{fwd:backward-energy}.
\end{proof}

The internal term is dissipative forward and antidissipative backward.
Neither sign is a sign theorem for the complete derivative: the signed
exterior term can reinforce or cancel it.

\begin{corollary}[Conditional forward force contraction]
\label{fwd:force-contraction}
Suppose an independent spatial enclosure gives
$x_N(t)-x_1(t)\le L_{\max}$ and $J(t)\ge\lVert Ph(t)\rVert_2$ on the
whole interval. With $\gamma=2N/L_{\max}^2$,
\begin{equation}
 r_U(t)\le \e^{-\gamma t}r_U(0)
       +\int_0^t\e^{-\gamma(t-v)}J(v)\dd v.
 \label{fwd:contraction}
\end{equation}
\end{corollary}
\begin{proof}
Every internal weight is at least $L_{\max}^{-2}$, and for a mean-zero
vector $\sum_{i<j}(a_i-a_j)^2=N\lVert a\rVert_2^2$. Thus
\eqref{fwd:energy} gives
$(r_U)_t\le-\gamma r_U+J$ wherever $r_U>0$.
At $r_U=0$ the upper right derivative of the norm is at most $\lVert Ph\rVert_2$,
so the same comparison holds in the upper-Dini sense. An integrating
factor proves \eqref{fwd:contraction}.
\end{proof}

For a constant $J$, the upper envelope is
$\e^{-\gamma t}r_U(0)+(J/\gamma)(1-\e^{-\gamma t})$.
It improves the initial norm only if its source floor is sufficiently
small. Without an acquired $J$, internal dissipation alone provides no
canonical contraction rate. The corresponding backward unsigned bound is
\begin{equation}
 (r_U)_u\le2\lambda_{\max}(L)r_U+\lVert Pf\rVert_2.
 \label{fwd:backward-norm}
\end{equation}
Replacing the denominators by one lower gap stop can make this scalar
allowance large; the fixed compact screen in
Appendix~\ref{app:compact-example} is scoped to that replacement and an
unsigned exterior norm.

Force dispersion is a velocity statistic. Indeed,
\begin{equation}
 g_i(t)-g_i(0)=2\int_0^t(U_{i+1}-U_i)\dd v,
 \qquad |U_{i+1}-U_i|\le\sqrt2\,r_U.
 \label{fwd:gap-integral}
\end{equation}
A small force dispersion need not imply a uniform gap profile: the
exterior may balance a nonuniform configuration. Direct common-source
node evolution retains the signed differences in
\eqref{fwd:gap-integral}; given the same node-force path information, the
norm replacement cannot improve that directional bound. No acquisition
cost advantage is asserted.

\subsection{A forward lower bound for the actual block minimum}

Assume the real segment connecting the two endpoint nodes lies inside
the independently cleared domain. Put
\begin{equation}
 g_i=x_{i+1}-x_i,\qquad D_i=g_i^2/4,\qquad
 q(t)=\min_{1\le i<N}D_i(t).
 \label{fwd:minimum}
\end{equation}
The minimum includes \emph{all} internal edges, including both end edges.
Split the complete interaction into the other real nodes in this block
and the complete exterior. Comparing \eqref{fwd:deflation} at two adjacent
nodes with the gap law gives
\begin{equation}
 M_i=M_{{\rm int},i}+M_{{\rm ext},i},\qquad
 M_{{\rm ext},i}=-\frac{R(t,x_{i+1})-R(t,x_i)}{g_i}
       =-\int_0^1R_z(t,x_i+\theta g_i)\dd\theta.
 \label{fwd:exterior}
\end{equation}
This is a complete signed exterior contribution. In particular,
$|R_z|\le C$ on every connecting segment implies
$M_{{\rm ext},i}\le C$, without exterior reality.

\begin{theorem}[Forward transport of a finite-block minimum]
\label{fwd:minimum-transport}
Suppose $M_{{\rm ext},i}\le C$ for every internal edge throughout the
forward interval, where $C\ge0$. Then, with
\begin{equation}
 a_N=\min_{1\le i<N}\left(\frac1i+\frac1{N-i}\right)
     =\frac{N}{\lfloor N/2\rfloor\lceil N/2\rceil},
 \label{fwd:finite-constant}
\end{equation}
the minimum obeys almost everywhere
\begin{equation}
 q_t\ge a_N-4Cq.
 \label{fwd:minimum-law}
\end{equation}
Consequently
\begin{equation}
 q(t)\ge\Phi_t(q(0)),\qquad
 \Phi_t(d)=
 \begin{cases}
 \e^{-4Ct}d+\dfrac{a_N}{4C}(1-\e^{-4Ct}),&C>0,\\[5pt]
 d+a_Nt,&C=0.
 \end{cases}
 \label{fwd:minimum-bound}
\end{equation}
\end{theorem}
\begin{proof}
At an active minimizing edge $i$, every block gap is at least $g_i$.
The $m$th retained node to the left or right of this pair is therefore
at distance at least $mg_i$ from its nearer endpoint. The positive
real-node terms in the interaction satisfy
\begin{align}
 M_{{\rm int},i}
 &\le\frac1{g_i^2}\left(
       \sum_{m=1}^{i-1}\frac1{m(m+1)}
       +\sum_{m=1}^{N-1-i}\frac1{m(m+1)}\right)\notag\\
 &=\frac1{g_i^2}\left(2-\frac1i-\frac1{N-i}\right).
 \label{fwd:internal-bound}
\end{align}
The empty sums at the end edges have value zero, so this also covers
$N=2$. Since $4D_i=g_i^2$, the physical pair law gives at an active edge
\[
 (D_i)_t=2-g_i^2M_i
 \ge\frac1i+\frac1{N-i}-4CD_i.
\]
Each $D_i$ is continuously differentiable, and the minimum of finitely
many such functions is locally Lipschitz. At almost every time its
derivative is the derivative of an active edge. The preceding bound
holds for every active edge, including ties, so it gives
\eqref{fwd:minimum-law} without a fixed minimizing label. Multiplication
by $\e^{4Ct}$ and integration proves \eqref{fwd:minimum-bound}.
The expression in \eqref{fwd:finite-constant} is minimized at the central
integer or integers, which gives the stated constant.
\end{proof}

For an acquired lower anchor $q(0)\ge d_0>0$, the scalar lower bound
$\Phi_t(d_0)$ improves $d_0$ at every $t>0$ precisely when
$C<a_N/(4d_0)$, with the stated continuous interpretation at $C=0$.
This is a test of a proposed source bound, not its acquisition. The
positive right side of \eqref{fwd:minimum-bound} can alternatively be used
in a first-lower-exit argument on an independently acquired outer domain.

The choice of minimum is essential. If only selected interior edges
were minimized, an excluded boundary gap could be smaller than the active
edge, invalidating the distance bound used in
\eqref{fwd:internal-bound}. A smaller observational window does not
automatically give its buffer to the outermost nodes of a larger census.
One must use the complete block minimum or pay a different boundary
comparison.

\begin{corollary}[Empirical lower-tail transport]
\label{fwd:empirical}
For finitely many blocks with the same $N$ and a common paid $C$, sample
the block index $X$ uniformly. For $C>0$,
\begin{equation}
 \mathbb P\{q_X(t)\le\xi\}
 \le\mathbb P\left\{q_X(0)\le
       \e^{4Ct}\xi-\frac{a_N}{4C}(\e^{4Ct}-1)\right\}.
 \label{fwd:tail}
\end{equation}
For $C=0$ the threshold on the right is $\xi-a_Nt$.
\end{corollary}
\begin{proof}
The increasing map $\Phi_t$ in \eqref{fwd:minimum-bound} gives the
pointwise event inclusion
$\{q_X(t)\le\xi\}\subseteq\{q_X(0)\le\Phi_t^{-1}(\xi)\}$.
Counting the finite block indices proves the claim. No independence or
random-matrix law is used.
\end{proof}

Different block sizes or exterior bounds use their respective maps
$\Phi_t$. A global zero-exception conclusion additionally needs coverage
of every eligible pair, including pairs crossing block boundaries, and
the finite integer threshold of Section~\ref{sec:statistics}. A small
empirical bad-block rate is not zero by itself; no varying-height or
all-height law is inferred from this finite statement.

\subsection{Complete-source acquisition of the exterior derivative}

The exterior estimate in Theorem~\ref{fwd:minimum-transport} can be
formulated on fixed trial locations, without assuming that the roots
remain positive-distance apart. For $w$ inside $\Gamma$,
\begin{equation}
 R(t,w)=\frac1{2\pi\ii}\oint_\Gamma\frac{B(t,z)}{z-w}\dd z,
 \qquad
 R_z(t,w)=\frac1{2\pi\ii}\oint_\Gamma
                         \frac{B(t,z)}{(z-w)^2}\dd z.
 \label{fwd:source-traces}
\end{equation}
To prove these formulas, split $B$ into $R$ and its finite set of
interior poles, with multiplicities. Cauchy's formula applies to $R$.
For each interior pole $\rho$, the rational functions
$[(z-\rho)(z-w)]^{-1}$ and
$[(z-\rho)(z-w)^2]^{-1}$ have all poles inside and decay at least as
$z^{-2}$ at infinity, so their contour integrals vanish. This remains
true at $w=\rho$ by the same rational argument or by continuation.
The count and full source determine the traces even if the interior
configuration is temporarily multiple or nonreal.

Suppose a circle of radius $r_\Gamma$ is used and every trial point $w$
is at distance at least $\delta>0$ from it. A uniform complete-source
error $e_B$ on the circle contributes at most
\begin{equation}
 E_R\le r_\Gamma e_B/\delta,\qquad
 E_{R_z}\le r_\Gamma e_B/\delta^2.
 \label{fwd:trace-error}
\end{equation}
These follow directly from contour length and the two kernels.
The quotient estimate \eqref{cmp:quotient-error} supplies $e_B$ only
after field, first-jet and positive contour-denominator errors are paid.
One must add quadrature, arithmetic, trial-location and genuine time-cell
errors; the angular-strip estimate \eqref{cmp:trapezoid} applies only
when the full integrand is cleared and bounded on that strip. A bound
for $|R_z|$ must hold on the whole trial segment containing all possible
connecting edges, not just at sampled nodes. The complete source error
already pays the exterior; no omitted zero tail is then silently set
to zero. No such favorable time-wide $C$ is acquired by these identities.

The force statistic has additional requirements. The fixed-coordinate
time derivative of the first trace in \eqref{fwd:source-traces}, when
its time-uniform source differentiation is justified, gives $R_t$ from
$B_t$ on the same contour. Alternatively one can evaluate the complete
source jets at simple roots. With a fixed scalar normalization $G$ of
$H$, differentiation along a root gives
\begin{equation}
 U_i=\frac{G_{zz}}{2G_z}(t,x_i),\qquad
 (U_i)_t=-\frac{G_{zzzz}}{2G_z}
       +\frac{G_{zz}G_{zzz}}{G_z^2}
       -\frac{G_{zz}^3}{2G_z^3}\quad\hbox{at }x_i.
 \label{fwd:source-jets}
\end{equation}
Indeed $\dd G_z(t,x_i)/\dd t=-G_{zzz}+G_{zz}^2/G_z$ and
$\dd G_{zz}(t,x_i)/\dd t=-G_{zzzz}+G_{zzz}G_{zz}/G_z$;
the quotient rule proves the second formula.
This route requires nonzero $G_z$ denominators and root-location errors.
For approximate node jets $h_1,h_2$ with errors $e_1,e_2$ and
$|h_1|\ge b_1>e_1$,
\begin{equation}
 \left|U-\frac{h_2}{2h_1}\right|
 \le\frac{e_2}{2(b_1-e_1)}
       +\frac{|h_2|e_1}{2b_1(b_1-e_1)}.
 \label{fwd:node-error}
\end{equation}
Higher quotient products in \eqref{fwd:source-jets} require their own
joint enclosure. If two exact and approximate positive node distances
are both at least $d_*$ and differ by at most $e$, their inverse-square
weights differ by at most $2e/d_*^3$. That error also enters a computed
Laplacian. Neither the common field nor mean subtraction removes these
reciprocal and boundary payments.

\subsection{What the forward estimate can pay on returning backward}

The internal estimate \eqref{fwd:internal-bound} is an \emph{upper}
interaction bound. Backward gap protection needs a lower interaction or
forcing bound, so it cannot reuse that inequality with its sign reversed.
For example, a two-sided exterior bound $|M_{{\rm ext},i}|\le C$ gives
only the coarse lower bound $M_i\ge-C$. On a paid backward interval it
therefore gives $D_u\ge-2-4CD$. If this is composed with
\eqref{fwd:minimum-bound} over equal forward and backward durations $s$,
and $z=\e^{-4Cs}$, the resulting scalar allowance is
\begin{equation}
 \mathcal Q_{\rm return}:=z^2q(0)
       +\frac{(1-z)(a_Nz-2)}{4C},\qquad C>0.
 \label{fwd:coarse-return}
\end{equation}
Here $\mathcal Q_{\rm return}$ is the lower allowance supplied by this
composition, not the actual return, which retraces the same trajectory.
To obtain \eqref{fwd:coarse-return}, the backward integrating factor
gives $D(s)\ge zD(0)-(1-z)/(2C)$; substitute the forward lower
anchor $zq(0)+a_N(1-z)/(4C)$. Since $a_N\le2$, the displayed right
side is at most $q(0)$. At $C=0$ the allowance is
$q(0)+(a_N-2)s$, with the same conclusion. Thus this particular pair
of scalar allowances does not turn a forward excursion into a better
returned lower bound.

The sharper comparison for the unchanged exact-initial-interaction
horizon is proved in Theorem~\ref{rf:horizon}. It likewise separates
an improved positive-time datum from an improved negative endpoint.
A useful forward acquisition must preserve additional joint information
for the backward consumer, such as the complete time-dependent forcing,
and give the direct source or regular $D/M$ comparison the same data.
The local forward theorem legitimately applies to noncanonical real
blocks as well; such controls may collide at an earlier backward time.
Its generic validity and the exact theta source identity are not, by
themselves, a canonical backward forcing estimate.
